% =============================================================================
% Lecture 5: Collective Communication
% Scaling Machine Learning Systems (Volume II), Chapter 5
% =============================================================================
% Thesis: a collective is a contract on what each rank holds afterward, and
% its time is the alpha-beta price paid once per round on a chunk.
% Running case: one gradient synchronization on DGX H100 nodes (NVLink
% 450 GB/s per direction in the node, one 50 GB/s InfiniBand NDR port per
% GPU, alpha = 1.5 us); the chapter's 1 GB reference buffer on 64 GPUs
% (8 nodes of 8) is the scoreboard, and the 70B model's 140 GB BF16
% gradient is the large case.
% Hook and synthesis: Horovod at Uber (War Story 5.1; Sergeev and Del Balso
% 2018, measured on 128 Pascal GPUs).
%
% Build: make deck V=2 D=05_collective_communication (figures: images/make_figs.py)
% Plan:  lecture_map.yml in this folder.
% Numbers: the rendered chapter; derived values and readings say so.
% =============================================================================

\documentclass[aspectratio=169, 12pt]{beamer}
\usepackage{../../assets/beamerthememlsys}
\mlsyssetup{
  volume={Volume II},
  chapter={Chapter 5},
  chaptertitle={Collective Communication},
  layer=distributed,
  logo={../../assets/img/logo-mlsysbook.png},
  instlogo={../../assets/img/logo-harvard.png},
}
\usepackage[T1]{fontenc}
\usepackage[scaled=0.9]{helvet}
\usepackage{courier}
\renewcommand{\familydefault}{\sfdefault}
\usepackage{booktabs,amsmath}
\usetikzlibrary{arrows.meta,positioning,calc}
\graphicspath{{images/}}
\setcounter{mlsystotalsections}{6}

\title{Collective Communication}
\author{Vijay Janapa Reddi}
\institute{Harvard University}
\date{}

% A large answer choice for predict and concept-question slides
\newcommand{\choice}[2]{\textbf{\textcolor{mlsysaccent}{#1}}\enspace #2}
\definecolor{commamber}{HTML}{B45309}
\definecolor{compblue}{HTML}{1F407A}

% Golden thread (D4): the Fleet Law term and C3 axis this lecture acts on, once, on the goals slide
\newcommand{\fleetterm}[2]{\par\vspace{0.08cm}{\small\textcolor{mlsysaccent}{\textbf{Fleet Law term:}}~#1\qquad\textcolor{mlsysaccent}{\textbf{C$^3$ axis:}}~#2\par}}

\begin{document}

% =============================================================================
% TITLE
% =============================================================================
\mlsysnote{[0.5 min] [CORE] Read the thesis aloud once. Last lecture fed and saved the fleet from storage; today the accelerators exchange data with each other, every training step. The notes total 70 minutes, of which 58 are core. FLEX: notes tagged [OPTIONAL] can be cut; cutting all of them leaves a complete lecture of 58 minutes.}
\mlsystitle{Collective Communication}{A collective is a contract on what each rank holds afterward, and its time is the $\alpha$-$\beta$ price paid once per round on a chunk.}{cover_collective_communication_blueprint.png}

% =============================================================================
% SECTION: TRAFFIC
% =============================================================================
\mlsysnote{[0.5 min] [CORE] Open on a real system that stopped scaling, then name the traffic each way of splitting a model creates.}
\mlsysdivider{Which data must cross between accelerators, and who must agree on it?}{Collective Communication, \S5.1, \S5.3, and \S5.4}
\section{Traffic}

\begin{frame}{Collectives live on the Distributed Execution layer}
  \centering
  \mlsyskitfig[height=5.4cm,keepaspectratio]{stack_distributed.pdf}
  \note{[0.5 min] [CORE] NARRATE: the fleet stack from facility at the bottom to governance at the top. Lectures 2 to 4 built the bottom two rows: accelerators, the fabric, and storage. Today moves up one row, to Distributed Execution, where a training job's logical split becomes physical traffic on that fabric. LINK: Lecture 1's Fleet Law charges this traffic to its communication term and credits whatever is hidden to its overlap term; every mechanism today shortens one of those two. FLEX: [CORE].}
\end{frame}

\begin{frame}{How much of 128 GPUs did TensorFlow use?}
  \mlsysactivity{Predict}{2 min}
  \centering
  \only<1|handout:0>{\safeimg[width=0.88\textwidth,height=3.9cm,keepaspectratio]{horovod_step1.pdf}}%
  \only<2-|handout:1>{\safeimg[width=0.88\textwidth,height=3.9cm,keepaspectratio]{horovod_step2.pdf}}\\[0.1cm]
  {\small Uber, 2017: a central parameter server sums every gradient. 128 Pascal GPUs, 25 Gb Ethernet.}\\[0.1cm]
  {\large\choice{A}{90\%}\qquad\choice{B}{70\%}\qquad\choice{C}{50\%}\qquad\choice{D}{25\%}}
  \mlsyscredit{Sergeev and Del Balso 2018 (Horovod), Figure 6, redrawn; bar heights read from the figure. Collective Communication, War Story 5.1.}
  \note{[2.5 min] [CORE] Step 1: dashed outlines are ideal throughput, the one-GPU rate times 128, for two image models. Ask students to vote on the share of ideal that standard distributed TensorFlow reached. Step 2: about half. Read from the paper's figure, Inception V3 reached about 53 percent and ResNet-101 about 41 percent; the paper's own words are that Uber lost about half its resources to communication overhead at 128 GPUs. TensorFlow's stock scheme was a parameter server, in which every worker sent its gradient to a central tier, which summed and sent back. ENGAGE: the GPUs are fast and the models fit. Where did half the hardware go? (Expected: waiting on gradient exchange.) WARN: students blame the network speed. Hold that thought; the next slide keeps the same network. FLEX: [CORE].}
\end{frame}

\begin{frame}{Horovod reached 88\% on the same GPUs and network}
  \begin{columns}[c]
    \begin{column}{0.45\textwidth}
      \centering\safeimg[width=\textwidth,height=5.3cm,keepaspectratio]{horovod_col.pdf}
    \end{column}
    \begin{column}{0.53\textwidth}
      \footnotesize\ttfamily
      \begin{tabular}{@{}l@{}}
        import horovod.tensorflow as hvd\\
        hvd.init()\\
        ...\\
        opt = tf.train.AdagradOptimizer(0.01)\\
        \textcolor{commamber}{\textbf{opt = hvd.DistributedOptimizer(opt)}}\\
        hvd.BroadcastGlobalVariablesHook(0)
      \end{tabular}\\[0.3cm]
      \normalfont\small Horovod replaced the parameter server with a ring AllReduce.
    \end{column}
  \end{columns}
  \mlsyscredit{Sergeev and Del Balso 2018: Figure 6 (bars read from the figure; 25GbE TCP), Section 8 (88 percent), Listing 1 (abridged).}
  \note{[2 min] [CORE] Same 128 GPUs, same 25 Gb Ethernet over TCP. Horovod, Uber's open-source library, reached the 88 percent efficiency mark on both models, about twice as fast as standard distributed TensorFlow. The code on the right is the paper's own listing, abridged: one wrapper around the optimizer, plus a broadcast of the initial weights from rank 0. Inside the wrapper, gradients are combined with ring AllReduce, the algorithm Baidu had popularized for deep learning earlier in 2017. NARRATE: only the communication schedule changed. The lecture's driving question follows. ENGAGE: if the bytes per GPU were the same, what else could differ? (Expected: where the bytes concentrate and how many rounds they take; both are today's topics.) FLEX: [CORE]. The synthesis returns to three more Horovod measurements.}
\end{frame}

\begin{frame}{Recall: what does one message cost?}
  \mlsysactivity{Retrieval}{1.5 min}
  \centering
  {\Large\bfseries Close your notes.}\\[0.3cm]
  {\normalsize Write the time to send one message of $n$ bytes over one link,\\and the per-GPU bandwidth inside a node and across nodes.}\\[0.5cm]
  \uncover<2->{%
    {\Large $T(n) = \alpha + n/\beta$}\\[0.25cm]
    {\normalsize $\alpha$: startup time per message\quad $\beta$: link bandwidth\\[0.1cm]
    DGX H100 GPU: NVLink 450 GB/s, InfiniBand port 50 GB/s (one direction)}}
  \mlsyscredit{Collective Communication, \S5.5 and \S5.7.1.}
  \note{[1.5 min] [CORE] Give 60 seconds. Step 2: Lecture 3's alpha-beta model. Alpha is the fixed startup cost of a message (software, PCIe, switching), about 1.5 microseconds on InfiniBand NDR; beta is the sustained bandwidth in one direction, 50 GB/s per port. LINK: Lecture 2's bandwidth hierarchy put NVLink at 450 GB/s per direction inside a DGX H100 node, about nine times one InfiniBand port. Lecture 4 moved bytes between storage and accelerators; today every byte moves between accelerators, every step. WARN: in this course alpha and beta always mean latency and bandwidth. FLEX: [CORE].}
\end{frame}

\begin{frame}{How long do N accelerators take to agree on a gradient?}
  \normalsize
  By the end of today, you will be able to:\\[0.05cm]
  \begin{itemize}\setlength{\itemsep}{0.04cm}
    \item Name the collective each parallelism split needs, by what each rank (one GPU) must hold.
    \item Price a collective round by round, and say whether a chunk is latency- or bandwidth-bound.
    \item Choose ring, tree, or hierarchical AllReduce for a message size and topology.
    \item Decide when compression or overlap shortens a step, and what stays exposed.
  \end{itemize}
  {\footnotesize\textcolor{midgray}{Running case: one gradient synchronization on DGX H100 nodes;\\the 1 GB reference buffer on 64 GPUs (8 nodes of 8); a 70B model's 140 GB BF16 gradient.}}
  \fleetterm{$T_{\text{comm}}(N)$ and $T_{\text{overlap}}$}{Communication ($C_2$)}
  \note{[1 min] [CORE] The title is the driving question. Each goal is one or two segments. The running case: DGX H100 nodes, eight GPUs each, NVLink inside the node, one 50 GB/s InfiniBand port per GPU. The chapter's 1 GB reference buffer, one large fused reduction, carries a scoreboard that returns after each segment; a 70-billion-parameter model's full gradient in BF16 (2 bytes per parameter, 140 GB) is the large case. WARN: the 140 GB case is communication arithmetic only; such a model's training state does not fit on one GPU, and Lecture 6 splits it. The accent line at the bottom names the Fleet Law term and C$^3$ axis that today's mechanisms act on; say it once. FLEX: [CORE].}
\end{frame}

\begin{frame}{Each way of splitting a model creates its own traffic}
  \centering
  \mlsyssteps[width=0.92\textwidth,height=4.4cm,keepaspectratio]{traffic_step}{3}\\[0.05cm]
  {\footnotesize\ttfamily\begin{tabular}{@{}l@{}}
  \textcolor{midgray}{\# torch.distributed}\enspace all\_reduce\enspace all\_gather\_into\_tensor\\
  reduce\_scatter\_tensor\enspace send\enspace recv\enspace all\_to\_all\_single
  \end{tabular}}
  \mlsyscredit{Collective Communication, table 1 (\S5.3), abridged. Function names: PyTorch v2.4.0, torch/distributed/distributed\_c10d.py.}
  \note{[2.5 min] [CORE] The map from strategy to traffic; Lecture 6 designs the strategies, today prices their traffic. Step 1: five splits. Data parallel copies the model and splits the batch, so replicas must agree on one gradient. Sharded data parallel (FSDP, ZeRO-3) also splits weights, gradients, and optimizer state across the replicas. Tensor parallel splits each layer's matrices inside a node. Pipeline gives each group of layers to a stage. Expert parallel sends each token to the GPU that holds its expert, as do recommendation models with sharded embedding tables. Step 2: the primitive each needs. Step 3: the term that binds; the narrowest cut is Lecture 3's bisection, the bandwidth across the thinnest slice that splits the fabric in two. Large gradients bind on bandwidth; pipeline messages are small and sit on the critical path, so latency binds; AllToAll binds on skew and on that cut. The monospaced line names PyTorch's own functions for these primitives. ENGAGE: which row repeats every step on the full gradient? (Expected: data parallel.) FLEX: [CORE].}
\end{frame}

\begin{frame}{A collective names what each rank holds afterward}
  \centering
  \mlsyssteps[width=0.95\textwidth,height=5.6cm,keepaspectratio]{contracts_step}{4}
  \mlsyscredit{Collective Communication, \S5.4 and figure 2. Ranks R0 to R3; color marks the source rank; numbers in AllToAll cells are chunk indices.}
  \note{[2.5 min] [CORE] Each panel shows four ranks before and after. Step 1, AllReduce: every rank contributes a full tensor and every rank ends with the same sum. Gradient synchronization is this case, in which every replica receives the reduced gradient, so all apply the same update. Step 2, AllGather: each rank holds one shard; afterward every rank holds all shards. Step 3, ReduceScatter: every rank contributes a full tensor; afterward rank $i$ holds only the $i$-th part of the sum. Step 4, AllToAll: rank $i$ sends its $j$-th chunk to rank $j$, a distributed transpose. Broadcast and Reduce, one root to all and all to one root, complete the set of six. ENGAGE: which two panels, run in order, give the first? (Expected: ReduceScatter then AllGather.) The contract fixes the result; ring, tree, and hierarchy are schedules for it. WARN: a collective is not a barrier. Ranks may enter at different times; a late rank delays only the ranks that need its data. FLEX: [CORE].}
\end{frame}

\begin{frame}{AllReduce is a ReduceScatter followed by an AllGather}
  \centering
  {\large AllReduce $=$ ReduceScatter $+$ AllGather}\\[0.15cm]
  {\footnotesize\ttfamily\begin{tabular}{@{}l@{}}
  class ShardingStrategy(Enum):\quad\textcolor{midgray}{\# torch/distributed/fsdp/api.py}\\
  \hspace*{1em}\textcolor{commamber}{\textbf{FULL\_SHARD = auto()}}\quad SHARD\_GRAD\_OP = auto()\\
  \hspace*{1em}NO\_SHARD = auto()\quad HYBRID\_SHARD = auto()
  \end{tabular}}\\[0.2cm]
  \small
  \begin{tabular}{@{}lccc@{}}
    \toprule
    Per step, $L$ layers & Calls & Bytes sent per rank & Startup, $\alpha_{\text{call}}$ per call \\
    \midrule
    Data parallel, unbucketed & 1 AllReduce & $\approx 2M$ & $\alpha_{\text{call}}$ \\
    FSDP & $2L$ & $\approx 2M$ & $2L\,\alpha_{\text{call}}$ \\
    FSDP, reshard after forward & $3L$ & $\approx 3M$ & $3L\,\alpha_{\text{call}}$ \\
    \bottomrule
  \end{tabular}\\[0.3cm]
  {\normalsize Sharding barely changes the bytes; it multiplies the calls.}
  \mlsyscredit{Collective Communication, \S5.4.1 and \S5.4.2 ($M$: model-sized payload, large group, equal widths). Code: PyTorch v2.4.0, lines 65 to 68.}
  \note{[2 min] [OPTIONAL] The two halves of AllReduce are exactly the two calls sharded data parallelism needs. FSDP keeps $1/N$ of each layer's weights per rank, so before a layer runs it calls AllGather to assemble the full weights, and after backward it calls ReduceScatter so each rank keeps only the gradient for its own shard. The code line is PyTorch's own enum: FULL\_SHARD shards weights, gradients, and optimizer state; NO\_SHARD falls back to an all-reduce like DDP; HYBRID\_SHARD shards inside a node and replicates across nodes. One AllReduce becomes $2L$ calls, or $3L$ if weights are freed after forward and gathered again for backward. Bytes per rank barely move, about $2M$ or $1.5\times$ that; the call count is what grows, and every call pays a startup cost. For a 32-layer model, 96 calls per step (derived). LINK: Lecture 6 explains why FSDP saves memory; today shows what it costs on the wire. WARN: calling a full AllReduce when each rank updates only its shard sends twice the bytes the step needs. FLEX: cut if short; the vote slide covers the primitives.}
\end{frame}

\begin{frame}{Which primitive does each exchange need?}
  \mlsysactivity{Vote}{2 min}
  \small
  \renewcommand{\arraystretch}{1.4}
  \begin{tabular}{@{}>{\raggedright\arraybackslash}m{0.5\textwidth}>{\raggedright\arraybackslash}m{0.44\textwidth}@{}}
    \toprule
    \textbf{Exchange} & \uncover<2->{\textbf{Primitive}} \\
    \midrule
    Every replica needs the summed gradient. & \uncover<2->{\textbf{AllReduce}} \\
    Each rank updates only its own shard of the weights. & \uncover<2->{\textbf{ReduceScatter}} \\
    A layer needs its full weights before it runs. & \uncover<2->{\textbf{AllGather}} \\
    Each token goes to the GPU that holds its expert. & \uncover<2->{\textbf{AllToAll}} \\
    \bottomrule
  \end{tabular}
  \mlsyscredit{Collective Communication, \S5.3 and \S5.4.1.}
  \note{[2 min] [CORE] Students write one primitive per row, compare with a neighbor, then reveal. Step 2: AllReduce for the replicated gradient; ReduceScatter when each rank keeps only its shard of the sum; AllGather to assemble sharded weights; AllToAll for destination-specific tokens. The rule is to choose the narrowest contract whose output the next computation consumes. ENGAGE: a student proposes AllReduce for the second row. What does it waste? (Expected: half the bytes; it assembles a full sum each rank immediately discards.) WARN: students reach for AllReduce for MoE routing. AllReduce sums matching elements; it cannot send different data to different ranks. FLEX: [CORE].}
\end{frame}

% =============================================================================
% SECTION: PRICING
% =============================================================================
\mlsysnote{[0.5 min] [CORE] We know what must move. Now put a price on each round.}
\mlsysdivider{What does one collective cost, and which term dominates?}{Collective Communication, \S5.5, \S5.6.1 to \S5.6.3}
\section{Pricing}

\mlsysnote{[1 min] [CORE] The equation of the lecture. Read it term by term. A ring has $2(N-1)$ rounds; each round pays one startup $\alpha$ and moves one chunk of $M/N$ bytes at bandwidth $\beta$. The first term grows with $N$. The second approaches $2M/\beta$ however large $N$ grows. Every algorithm today changes one of these: the number of rounds, the chunk size, or the $\beta$ of the link a round uses. LINK: this is Lecture 1's communication term for one gradient synchronization. FLEX: [CORE].}
\mlsysfocus{Ring AllReduce time}{%
$T_{\text{ring}} \;=\; \underbrace{2(N-1)\,\alpha}_{\text{one startup per round}} \;+\; \underbrace{2\,\frac{N-1}{N}\,\frac{M}{\beta}}_{\text{one chunk per round}}$\\[0.4cm]
{\normalsize $N$ ranks\quad $M$ bytes per rank\quad $\alpha$ startup\quad $\beta$ bandwidth}}

\begin{frame}{Each link has a size where startup and transfer cost the same}
  \centering
  \only<1|handout:0>{\safeimg[width=0.95\textwidth,height=5.6cm,keepaspectratio]{chunk_step1.pdf}}%
  \only<2-|handout:1>{\safeimg[width=0.95\textwidth,height=5.6cm,keepaspectratio]{chunk_step2.pdf}}
  \mlsyscredit{Collective Communication, table 4 (vendor latency and bandwidth ranges, midpoints) and Napkin Math 5.1.}
  \note{[2.5 min] [CORE] Step 1: the critical size $n^* = \alpha\beta$ for five interconnects. Below it, a message spends more time starting than transferring, so it is latency-bound; above it, bandwidth-bound. InfiniBand NDR: 1.5 microseconds times 50 GB/s is 75 KB. NVLink stays latency-bound up to about 0.7 MB because its bandwidth is so high. Step 2: the chunk size sets the regime. A 25 MB gradient bucket in a 64-rank ring travels in 390.6 KB chunks, above the NDR knee. On the 1,024-GPU reference cluster the chunk is 24.4 KB, below it. ENGAGE: on the 1,024-rank ring, does a faster link help much? (Expected: little; startup dominates each round.) WARN: students classify by tensor size. A 25 MB bucket is 333 times the knee, yet its ring steps are latency-bound at 1,024 ranks. FLEX: [CORE].}
\end{frame}

\begin{frame}{The same 25 MB bucket takes 3.5$\times$ longer on 1,024 GPUs}
  \centering
  \normalsize
  \begin{tabular}{@{}lrrrr@{}}
    \toprule
    Ring of a 25 MB bucket, NDR & Chunk & Rounds & Byte term & Total \\
    \midrule
    64 GPUs & 390.6 KB & 126 & 1.0 ms & 1.2 ms \\
    1,024 GPUs & \textbf{24.4 KB} & \textbf{2,046} & 1.0 ms & \textbf{4.1 ms} \\
    \bottomrule
  \end{tabular}\\[0.3cm]
  \uncover<2->{{\normalsize Latency-bound levers: \textbf{fusion}, \textbf{fewer rounds}, \textbf{hierarchy}.}\\[0.1cm]
  {\small Fusing 100 messages of 4 KB: $100\alpha = 150~\mu$s becomes $\alpha + 400~\text{KB}/\beta = 9.5~\mu$s.}}\\[0.25cm]
  \begin{minipage}{0.92\textwidth}\footnotesize\ttfamily
  \$ ./build/all\_reduce\_perf -b 8 -e 128M -f 2 -g 8\\
  \textcolor{midgray}{\# small sizes measure the constant overhead (nccl-tests docs)}
  \end{minipage}
  \mlsyscredit{Collective Communication, Napkin Math 5.1 ($\alpha = 1.5~\mu$s, $\beta = 50$ GB/s) and \S5.5. Ratio derived (calcs.py). Command: NVIDIA nccl-tests README; comment paraphrases doc/PERFORMANCE.md.}
  \note{[2 min] [CORE] Napkin Math 5.1. Ring time is $2(N-1)(\alpha + M/(N\beta))$. On 64 ranks, 126 rounds of 9.3 microseconds give 1.2 ms, of which 1.0 ms is bytes. On 1,024 ranks the byte term is still 1.0 ms, but 2,046 rounds each pay 1.5 microseconds of startup, and the total is 4.1 ms (derived ratio 3.5 from the unrounded times, calcs.py). Step 2: the levers for latency-bound steps. Fusion packs many small messages into one: 100 routing messages of 4 KB cost about 150 microseconds separately and about 9.5 together, roughly 16 times less. Horovod's Tensor Fusion and PyTorch's gradient buckets are this lever. Fewer rounds is the tree; hierarchy is the next segment. The command is NVIDIA's own benchmark, sweeping sizes from 8 bytes to 128 MB on 8 GPUs; its documentation says small sizes measure the fixed overhead, which is $\alpha$. ENGAGE: does compression help the 4 KB messages? (Expected: no; it shrinks bytes, and their cost is startup.) FLEX: [CORE].}
\end{frame}

\begin{frame}{A parameter server pushes 32.5$\times$ the ring's bytes through one port}
  \centering
  \only<1|handout:0>{\safeimg[width=0.92\textwidth,height=4.6cm,keepaspectratio]{scoreboard_step1.pdf}}%
  \only<2-|handout:1>{\safeimg[width=0.92\textwidth,height=4.6cm,keepaspectratio]{scoreboard_step3.pdf}}\\[0.15cm]
  {\small\uncover<2->{Some rank must send at least $2\frac{N-1}{N}M$ bytes; the ring meets that bound.}}
  \mlsyscredit{Collective Communication, \S5.6.1 and \S5.6.2. 1 GB buffer, 64 workers, 50 GB/s. Server time $N M/\beta$ derived; ring byte term 39.4 ms; tree 240 ms (\S5.6.4).}
  \note{[2 min] [CORE] This is the running scoreboard, with the chapter's 1 GB reference buffer on 64 GPUs at 50 GB/s. Step 1: a parameter server receives 64 GB and sends 64 GB through its one port, 64 GB over 50 GB/s, about 1,280 ms each way (derived). Step 2: the ring spreads the same reduction so every rank sends about $2M$, 39.4 ms. The ratio is 32.5 (derived), close to $N/2 = 32$. The bound has two halves. Combining $N$ contributions takes at least $N-1$ transfers, and delivering the result to $N-1$ other ranks takes $N-1$ more, so $2(N-1)M$ bytes in total, of which some rank sends at least its share. The simple tree, 240 ms, is the next segment. ENGAGE: total bytes are the same order for star and ring. Why is the star slow? (Expected: concentration in one port.) WARN: this is Uber's problem, and sharding the server spreads it only by adding servers. FLEX: [CORE].}
\end{frame}

\begin{frame}{Ring AllReduce takes six steps, one chunk per GPU each}
  \begin{columns}[c]
    \begin{column}{0.33\textwidth}
      \centering\safeimg[width=\textwidth,height=3.6cm,keepaspectratio]{ring.pdf}\\[0.1cm]
      {\footnotesize GPU 0 to 3 hold\\{[1,5,3,7]}, {[2,6,4,8]},\\{[3,7,5,9]}, {[4,8,6,10]}.\\Sum $=$ [10, 26, 18, 34].}
    \end{column}
    \begin{column}{0.64\textwidth}
      \centering\small
      \begin{tabular}{@{}lcccc@{}}
        \toprule
        Step & GPU 0 & GPU 1 & GPU 2 & GPU 3 \\
        \midrule
        \uncover<2->{ReduceScatter 1} & \uncover<2->{d: 17} & \uncover<2->{a: 3} & \uncover<2->{b: 13} & \uncover<2->{c: 11} \\
        \uncover<3->{ReduceScatter 2} & \uncover<3->{c: 14} & \uncover<3->{d: 25} & \uncover<3->{a: 6} & \uncover<3->{b: 21} \\
        \uncover<4->{ReduceScatter 3} & \uncover<4->{\textbf{b: 26}} & \uncover<4->{\textbf{c: 18}} & \uncover<4->{\textbf{d: 34}} & \uncover<4->{\textbf{a: 10}} \\
        \midrule
        \uncover<5->{AllGather 1} & \uncover<5->{a: 10} & \uncover<5->{b: 26} & \uncover<5->{c: 18} & \uncover<5->{d: 34} \\
        \uncover<6->{AllGather 2} & \uncover<6->{d: 34} & \uncover<6->{a: 10} & \uncover<6->{b: 26} & \uncover<6->{c: 18} \\
        \uncover<7->{AllGather 3} & \uncover<7->{c: 18} & \uncover<7->{d: 34} & \uncover<7->{a: 10} & \uncover<7->{b: 26} \\
        \bottomrule
      \end{tabular}\\[0.2cm]
      {\small Cell: the chunk each GPU received and the value it now holds.}
    \end{column}
  \end{columns}
  \mlsyscredit{Collective Communication, Example 5.1 and table 5. Chunks a to d are the four elements.}
  \note{[3.5 min] [CORE] Example 5.1. Each GPU splits its vector into four chunks, a to d. Step 1 (ReduceScatter 1): every GPU sends one running chunk to its right neighbor, which adds its own. GPU 0 receives GPU 3's d (10) and adds its own 7: 17. Steps 2 and 3: each partial sum moves on and grows. After step 3 each GPU owns one fully reduced chunk, in bold: GPU 0 holds b = 26. Steps 4 to 6 (AllGather): finished chunks circulate without arithmetic until every GPU holds [10, 26, 18, 34]. Point out three properties: $2(N-1) = 6$ steps; every GPU sends and receives in every step; each step moves one chunk, $M/N$. So each GPU sends $2(N-1)M/N$, the lower bound. ENGAGE: one GPU runs at half speed. What happens to the whole AllReduce? (Expected: every step waits for it; the ring runs at its slowest member.) FLEX: [CORE].}
\end{frame}

\begin{frame}{How long does one 140 GB AllReduce take?}
  \mlsysactivity{Estimate}{2.5 min}
  \begin{columns}[T]
    \begin{column}{0.4\textwidth}
      \begin{mlsyscard}{mlsysaccent}
        \textbf{Given}\\[0.1cm]
        70B parameters, BF16 gradient\\
        64 GPUs in one ring\\
        $\beta = 50$ GB/s, $\alpha = 1.5~\mu$s
      \end{mlsyscard}
    \end{column}
    \begin{column}{0.56\textwidth}
      \normalsize
      \uncover<2->{$M = 70\times10^9 \times 2~\text{B} = 140$ GB\\[0.15cm]}
      \uncover<3->{Bytes: $2\cdot\frac{63}{64}\cdot\frac{140}{50}$ s $= 5{,}512.5$ ms\\[0.15cm]}
      \uncover<4->{Startup: $126 \times 1.5~\mu$s $= 0.2$ ms\\[0.15cm]}
      \uncover<5->{\textbf{Total $\approx$ 5.5 s per step}, almost all bandwidth.\\[0.1cm]
      {\small Halving it saves about 4,900 GPU-hours over 100,000 steps.}}
    \end{column}
  \end{columns}
  \vspace{0.2cm}\centering
  {\footnotesize\ttfamily B = algbw * (2*(n-1)/n)\quad\textcolor{midgray}{\# nccl-tests bus bandwidth}}
  \mlsyscredit{Collective Communication, Napkin Math 5.2 (communication-only baseline; each rank reduces the full gradient). Formula: nccl-tests doc/PERFORMANCE.md.}
  \note{[2.5 min] [CORE] Give students 90 seconds with the card. Step 2: 140 GB of gradient. Step 3: the byte term, 1.96875 times 140 GB over 50 GB/s, 5,512.5 ms. Step 4: 126 rounds of 1.5 microseconds, 0.2 ms. Step 5: about 5.5 seconds per synchronization, and startup is 0.0034 percent (calcs.py). The bottom line is from NVIDIA's nccl-tests documentation, whose bus bandwidth multiplies measured algorithm bandwidth by the same $2(n-1)/n$, so a healthy ring reports close to the link's peak. If fully exposed, halving it saves 2.76 seconds per step; 2.76 s times 64 GPUs times 100,000 steps is about 4,900 GPU-hours (calcs.py). ENGAGE: which levers help here? (Expected: byte levers: fewer bytes per link, more links in parallel, overlap. Fusion does nothing.) WARN: this assumes every rank holds the full gradient; partitioned training reduces shards in smaller groups (Lecture 6). FLEX: [CORE].}
\end{frame}

% =============================================================================
% SECTION: ALGORITHMS
% =============================================================================
\mlsysnote{[0.5 min] [CORE] The ring is optimal in bytes and linear in rounds. Can we keep the bytes and cut the rounds, and use the fast links inside a node?}
\mlsysdivider{Which AllReduce schedule wins, and when?}{Collective Communication, \S5.6.4 to \S5.6.7 and \S5.7}
\section{Algorithms}

\begin{frame}{A tree needs 20 rounds where a ring needs 2,046}
  \centering
  \safeimg[width=0.92\textwidth,height=5.2cm,keepaspectratio]{rounds.pdf}\\[0.1cm]
  {\normalsize The simple tree pays in bytes, because the full message crosses every level.}
  \mlsyscredit{Collective Communication, \S5.6.4 and its round-count figure, redrawn.}
  \note{[2 min] [OPTIONAL] A tree reduces pairs up to a root in $\log_2 N$ rounds and broadcasts back down in $\log_2 N$ more. At 1,024 ranks: 20 rounds against the ring's 2,046, about 100 times fewer. The cost is bytes. In the simple tree each level forwards the whole message, so the byte term is $2\log_2 N \cdot M/\beta$. On the 64-GPU scoreboard that is 12 full transfers of 1 GB, 240 ms against the ring's 39.4 ms. ENGAGE: which payloads want the tree? (Expected: small ones: scalar reductions, control messages, small activations, where rounds dominate.) WARN: 'tree' names a family. Recursive halving-doubling and NCCL's double binary tree keep log rounds with ring-level bytes; the simple tree is the teaching baseline. FLEX: if short, cut this slide; the predict slide that follows states both round counts.}
\end{frame}

\begin{frame}{Ring or tree for a 1 MB buffer on 64 GPUs?}
  \mlsysactivity{Predict}{2.5 min}
  \begin{columns}[c]
    \begin{column}{0.36\textwidth}
      \begin{mlsyscard}{mlsysaccent}
        $N = 64$, $M = 1$ MB\\
        $\alpha = 10~\mu$s, $\beta = 10$ GB/s
      \end{mlsyscard}
      \vspace{0.15cm}
      \small
      \uncover<2->{%
      \begin{tabular}{@{}lrr@{}}
        \toprule
        $\mu$s & Ring & Tree \\
        \midrule
        Startup & 1,260 & 120 \\
        Bytes & 196.9 & 1,200 \\
        \textbf{Total} & 1,456.9 & \textbf{1,320} \\
        \bottomrule
      \end{tabular}}
    \end{column}
    \begin{column}{0.62\textwidth}
      \centering
      \uncover<3->{\safeimg[width=\textwidth,height=4.1cm,keepaspectratio]{crossover.pdf}\\[0.05cm]
      {\small $M^* \approx \dfrac{N\alpha\beta}{\log_2 N}$: 1.1 MB here, 7.7 MB on 1,024 NDR GPUs.}\\[0.05cm]
      {\footnotesize\ttfamily NCCL\_ALGO=Ring \textbar{} Tree \textbar{} NVLS \textbar{} ...}}
    \end{column}
  \end{columns}
  \mlsyscredit{Collective Communication, Napkin Math 5.3 and \S5.6.7. Curves: the chapter's two cost models. NCCL\_ALGO: NCCL User Guide, environment variables.}
  \note{[2.5 min] [CORE] Students vote ring or tree. Step 2: ring startup is 126 rounds of 10 microseconds, 1,260; tree startup is 12 rounds, 120. Ring bytes are 196.9 microseconds; tree bytes are 12 full transfers, 1,200. Tree wins by about 9 percent. Step 3: the two models against message size. Below the crossover the tree's saved rounds win; above it the ring's saved bytes win. Setting the two equal and keeping the large-$N$ terms gives $M^* \approx N\alpha\beta/\log_2 N$: about 1.1 MB here, and 7.7 MB on the 1,024-GPU NDR reference cluster. So a 25 MB bucket or a full gradient goes to the ring; a few hundred kilobytes go to the tree. ENGAGE: the crossover grows with $N$. Why? (Expected: the ring's round count grows linearly.) The bottom line is NCCL's real setting. By default the library picks the algorithm per call, and NCCL\_ALGO restricts it to Ring, Tree, NVLS (in-switch reduction over NVLink), and others. WARN: NCCL chooses among more than two schedules; the model names a candidate to benchmark, and an override stays only if measured step time improves. FLEX: [CORE].}
\end{frame}

\begin{frame}{Two schedules keep log rounds and ring-level bytes}
  \centering
  \small
  \renewcommand{\arraystretch}{1.25}
  \begin{tabular}{@{}lllp{4.3cm}@{}}
    \toprule
    Schedule & Startup rounds & Bytes sent per rank & Condition \\
    \midrule
    Ring & $2(N-1)$ & $2\frac{N-1}{N}M$ (optimal) & none \\
    Simple tree & $2\log_2 N$ & $2\log_2 N \cdot M$ & none \\
    Halving-doubling & $2\log_2 N$ & $2\frac{N-1}{N}M$ (optimal) & $N$ a power of two; distant partners \\
    Double binary tree & $\approx 2\log_2 N$ & $\approx 2M$ & two complementary trees, pipelined \\
    \bottomrule
  \end{tabular}\\[0.3cm]
  {\normalsize Fewer rounds bring a placement cost, because distant partners share the fabric.}
  \mlsyscredit{Collective Communication, table 6 and \S5.6.5 to \S5.6.6.}
  \note{[1 min] [OPTIONAL] Halving-doubling (butterfly) pairs rank $i$ with rank $i$ XOR $2^k$ and halves the region exchanged each round, so it reaches the byte bound in log rounds; but its last rounds pair distant ranks, which may sit on different racks and share the inter-rack fabric. NCCL's double binary tree runs two trees at once, each on half the data; ranks that are interior in one tree are leaves in the other, and each tree pipelines small chunks, so every rank sends about $2M$. ENGAGE: on 8 GPUs, 4 per node, which halving-doubling round crosses the node? (Expected: the third, partner distance 4.) FLEX: cut if short; the takeaway is that the libraries pick among these by message size, rank count, and topology.}
\end{frame}

\begin{frame}{Inside a node, NVLink carries 9$\times$ an InfiniBand port}
  \begin{columns}[c]
    \begin{column}{0.48\textwidth}
      \centering
      \safeimg[width=\textwidth,height=4.9cm,keepaspectratio]{nvidia-hgx-b200-board.jpg}\\[0.05cm]
      {\footnotesize An HGX B200 board: eight GPUs under heat sinks, joined by NVLink on the board.}
    \end{column}
    \begin{column}{0.48\textwidth}
      \centering
      \safeimg[width=\textwidth,height=4.0cm,keepaspectratio]{bw_tiers.pdf}\\[0.1cm]
      {\small DGX H100, per GPU. A flat ring runs at its slowest edge.}
    \end{column}
  \end{columns}
  \mlsyscredit{Photo: Pokiiri, CC BY-SA 4.0, Wikimedia Commons (HGX B200; the lecture's numbers use the H100 generation). Bandwidths: Collective Communication, \S5.7.1.}
  \note{[2 min] [OPTIONAL] Real, then abstract. The photo is an eight-GPU board of the generation after the H100; the eight heat sinks cover eight GPUs, and NVLink switches on the board join them. The running case uses DGX H100 numbers: 450 GB/s per GPU per direction over NVLink (the 900 GB/s on the datasheet counts both directions), and one 50 GB/s InfiniBand port per GPU to leave the node. LINK: Lecture 2's bandwidth hierarchy. Every ordering of a flat ring over 64 GPUs on 8 nodes has at least one edge per node that crosses InfiniBand, and every edge carries the same traffic each step, so the whole ring runs at 50 GB/s and takes about 40 ms for the 1 GB buffer. ENGAGE: which links sit idle during that flat ring? (Expected: most NVLink capacity, and seven of each node's eight ports.) FLEX: if short, say the 450 against 50 GB/s line on the next slide and skip the photo.}
\end{frame}

\begin{frame}{Hierarchical AllReduce turns 40 ms into 8.3 ms on 8 $\times$ 8 GPUs}
  \centering
  \mlsyssteps[width=0.95\textwidth,height=5.0cm,keepaspectratio]{hier_step}{3}\\[0.15cm]
  {\normalsize\uncover<3->{1.94 $+$ 4.40 $+$ 1.94 $=$ \textbf{8.3 ms}, against about 40 ms for a flat ring on one port per hop.}}
  \mlsyscredit{Collective Communication, \S5.7.1, figure 7, and Napkin Math 5.4 (1 GB, 8 nodes of 8, one 50 GB/s port per GPU). Drawn with 4 nodes of 4.}
  \note{[3 min] [CORE] Drawn with four nodes of four GPUs; the numbers are for eight of eight. Step 1: inside each node, a ReduceScatter over NVLink leaves each GPU one eighth of the node's partial sum. Each GPU sends 875 MB at 450 GB/s, about 1.94 ms. Step 2: GPUs in the same position on every node form a ring and AllReduce their 125 MB shard over InfiniBand, about 4.40 ms. The eight rings run at once, one per port. Step 3: an AllGather over NVLink rebuilds the full sum in each node, 1.94 ms. Total 8.3 ms against about 40 ms, 4.8 times faster. The scoreboard's last bar. On the 1,024-GPU reference cluster, Archetype A's 175B gradient (350 GB in BF16) drops from about 14.0 s to 3.1 s. ENGAGE: what if the node's eight GPUs shared one uplink? (Expected: the eight shard rings run one after another and most of the gain disappears.) FLEX: [CORE].}
\end{frame}

\begin{frame}{Rail alignment gives each shard ring its own NIC}
  \begin{columns}[c]
    \begin{column}{0.46\textwidth}
      \centering
      \safeimg[width=\textwidth,height=3.4cm,keepaspectratio]{infiniband-switch-cables.jpg}\\[0.05cm]
      {\footnotesize InfiniBand switch ports and cables.}
    \end{column}
    \begin{column}{0.50\textwidth}
      \small
      \textbf{Rail $k$:} GPU $k$ of every node, one switch.\\[0.2cm]
      Aligned ranks: 8 shard rings on 8 NICs.\\[0.2cm]
      \textcolor{RedLine}{Misplaced ranks: all 8 rings on one NIC, \textbf{1/8} the injection bandwidth.}\\[0.3cm]
      {\footnotesize Record rank, GPU, NIC, node, and rail before trusting a benchmark.}
    \end{column}
  \end{columns}
  \mlsyscredit{Photo: ChrisDag, CC BY 2.0, Wikimedia Commons. Rails: Collective Communication, \S5.7.4 and table 8.}
  \note{[1.5 min] [OPTIONAL] Step 2 of the hierarchical schedule depends on wiring. In a DGX cluster, GPU $k$ of every node connects through its own NIC to rail switch $k$, so the eight shard rings travel on eight independent rails. LINK: Lecture 3 introduced rail-optimized topology. If the scheduler numbers ranks without regard to the hardware, the shard groups and the rails drift apart; in the worst case all eight exchanges converge on one NIC, one eighth of the rail-aligned bandwidth. A common symptom is that nccl-tests shows full bandwidth while the training job gets a fraction. WARN: contiguous rank numbers on a node do not prove rail alignment. LINK: Lecture 8 places jobs on the topology. Two other ways to shrink payloads before the slow tier, torus reduction on TPU pods and in-switch reduction (SHARP), are in the reading, \S5.7.2 and \S5.7.3. FLEX: cut if short.}
\end{frame}

% =============================================================================
% SECTION: ALLTOALL
% =============================================================================
\mlsysnote{[0.5 min] [CORE] Hierarchy, trees, and rails were built for AllReduce's uniform payload. Expert routing sends different data to every rank.}
\mlsysdivider{Why can AllToAll stall before AllReduce does?}{Collective Communication, \S5.8}
\section{AllToAll}

\begin{frame}{MoE dispatch sends 64 tokens from every GPU to every other}
  \begin{columns}[c]
    \begin{column}{0.43\textwidth}
      \centering\safeimg[width=\textwidth,height=5.6cm,keepaspectratio]{a2a_matrix.pdf}
    \end{column}
    \begin{column}{0.55\textwidth}
      \small
      8 GPUs, 64 experts, 512 tokens each.\\
      Token: 2,048 values in BF16, 4.096 KB.\\[0.2cm]
      Sent per GPU: $7 \times 64 \times 4.096$ KB $=$ \textbf{1.84 MB}\\[0.2cm]
      \uncover<2->{Bound: $\geq 1.84$ MB $/$ 50 GB/s $=$ \textbf{36.7 $\mu$s}\\[0.15cm]
      {\footnotesize It ends at the largest of four terms: startup, injection, reception, and bytes over the bisection.}}\\[0.15cm]
      {\footnotesize\ttfamily\begin{tabular}{@{}l@{}}dist.all\_to\_all\_single(output, input,\\\hspace*{1em}output\_split\_sizes,\\\hspace*{1em}input\_split\_sizes)\end{tabular}}
    \end{column}
  \end{columns}
  \mlsyscredit{Collective Communication, Napkin Math 5.5 and Systems Perspective 5.1 (uniform routing, InfiniBand NDR). Code: PyTorch v2.4.0 distributed\_c10d.py.}
  \note{[2.5 min] [CORE] A mixture-of-experts layer with 64 experts on 8 GPUs, 8 experts per GPU. A gating network picks one expert per token. Uniform routing: each GPU keeps 64 of its 512 tokens and sends 64 to each of the 7 others; the matrix is the AllToAll's traffic, sender by receiver. That is 1.84 MB per GPU, and the results come back by a second AllToAll. Step 2: even if all seven sends start at once they share one 50 GB/s port, so dispatch takes at least 36.7 microseconds. The four-term bound: per-message startup, the sender's injection port, the receiver's port, and the bytes crossing the fabric's narrowest cut. $N(N-1)$ logical pairs grow quadratically, but per-rank bytes can stay flat. ENGAGE: seven concurrent sends at 50 GB/s each: 350 GB/s? (Expected: no; one port.) The code line is PyTorch's AllToAll; the split sizes let each destination receive a different count, the variable-size form skewed routing needs. LINK: Lecture 6 designs expert parallelism; today prices its exchange. FLEX: [CORE].}
\end{frame}

\begin{frame}{One rank gets 1.5$\times$ the tokens. What happens?}
  \mlsysactivity{Diagnose}{2.5 min}
  \centering
  {\normalsize Routing is data dependent. One GPU's experts receive 672 remote tokens, not 448.\\How long does dispatch take for the whole layer?}\\[0.3cm]
  \only<1|handout:0>{\safeimg[width=0.92\textwidth,height=2.9cm,keepaspectratio]{a2a_hot_step1.pdf}}%
  \only<2-|handout:1>{\safeimg[width=0.92\textwidth,height=2.9cm,keepaspectratio]{a2a_hot_step2.pdf}}\\[0.2cm]
  {\normalsize\uncover<2->{Every rank waits for the hottest receiver, \textbf{1.5$\times$ longer}.}}
  \mlsyscredit{Collective Communication, Napkin Math 5.6 (assumed skew on one rank; 50 GB/s receive port).}
  \note{[2.5 min] [CORE] Ask: the other seven GPUs finish their receives sooner. Does the layer? Step 2: no. The receive bound for the hot GPU is 672 tokens times 4.096 KB, 2.75 MB, over 50 GB/s: 55.1 microseconds against 36.7. Work that needs the hot expert's outputs waits, so the whole layer's dispatch rises by the same 1.5 times. Its expert computation may also grow, and becomes a straggler. Remedies follow the bound. Variable-size AllToAll avoids sending padded empty slots but not the hot receive. Hierarchical AllToAll coalesces small chunks by destination node to amortize startup, but cannot shrink bytes. Balancing the router attacks the skew itself. ENGAGE: a capacity factor of 1.25 reserves 25 percent more slots per expert. What does it cost if every slot is sent? (Expected: about 1.25 times the bytes.) FLEX: [CORE].}
\end{frame}

% =============================================================================
% SECTION: EXPOSED TIME
% =============================================================================
\mlsysnote{[0.5 min] [CORE] After the schedule is chosen, two levers remain: send fewer bytes, or hide the transfer under computation.}
\mlsysdivider{How much of the remaining communication must the step wait for?}{Collective Communication, \S5.9 to \S5.11}
\section{Exposed Time}

\begin{frame}{Compression pays when it removes more time than it adds}
  \centering
  \safeimg[width=0.92\textwidth,height=3.2cm,keepaspectratio]{payback.pdf}\\[0.3cm]
  {\normalsize Pays only if $t_{\text{codec}} < T_{\text{comm}}\,(1 - 1/r)$, with compression ratio $r$.}\\[0.15cm]
  {\small INT8: up to 4$\times$ vs FP32.\quad Top-1\% values + indices: up to 50$\times$.\quad 1-bit: up to 32$\times$.}\\[0.15cm]
  {\footnotesize\ttfamily ddp\_model.register\_comm\_hook(process\_group, fp16\_compress\_hook)}
  \mlsyscredit{Collective Communication, Napkin Math 5.7 (assumed 8$\times$ effective, 2 ms codec) and table 12. Code: PyTorch v2.4.0, ddp\_comm\_hooks/default\_hooks.py docstring.}
  \note{[1.5 min] [CORE] Compression shrinks only the byte term, so it helps payloads far above $\alpha\beta$, such as a full gradient. Napkin Math 5.7: a 40 ms synchronization (the scoreboard's flat ring) compressed with an assumed effective 8 times becomes 5 ms of transfer plus 2 ms of encoding, 7 ms, about 5.7 times faster. On NVLink, where the same reduction might take 1 ms, 2 ms of encoding makes it slower. The ratios on the bottom line are against FP32; the 70B gradient is already BF16, so each ratio halves. Top-k keeps the largest 1 percent of values but must send an index with each, which halves its headline ratio. The code line is PyTorch's own example, a DDP communication hook that casts each bucket to FP16 before the AllReduce, a 2$\times$ compression against FP32. ENGAGE: compress a 4 KB control message? (Expected: no; its cost is startup.) FLEX: [CORE].}
\end{frame}

\begin{frame}{Does compression help on a fast network?}
  \mlsysactivity{Predict}{2.5 min}
  \centering
  \only<1|handout:0>{\safeimg[width=0.88\textwidth,height=4.4cm,keepaspectratio]{runlevel_step1.pdf}}%
  \only<2|handout:0>{\safeimg[width=0.88\textwidth,height=4.4cm,keepaspectratio]{runlevel_step2.pdf}}%
  \only<3-|handout:1>{\safeimg[width=0.88\textwidth,height=4.4cm,keepaspectratio]{runlevel_step3.pdf}}\\[0.15cm]
  {\normalsize 100,000 steps, 300 ms compute each. Compression cuts communication 4$\times$\\but needs 10\% more steps to reach the same quality.}
  \mlsyscredit{Collective Communication, \S5.9.5 (serial compute plus communication; scenario inputs).}
  \note{[2.5 min] [CORE] Step 1: two networks. Ask students to predict, for each, whether compression shortens time to target. Step 2, slow network: 500 ms steps (300 compute, 200 communication) take 50,000 s. Compression cuts communication to 50 ms, but 110,000 steps at 350 ms take 38,500 s, 23 percent faster. Step 3, fast network: 330 ms steps take 33,000 s; compression saves 22.5 ms per step, but 110,000 steps at 307.5 ms take 33,825 s, 2.5 percent slower. ENGAGE: what decides the verdict? (Expected: the share of the step that communication takes, against the extra steps.) WARN: the chapter's pitfall is judging compression by communication speedup. Judge it by time to the same quality. FLEX: [CORE].}
\end{frame}

\begin{frame}{Error feedback carries what rounding dropped into the next step}
  \centering
  \small
  {\normalsize Compressor: send each value rounded to a whole number.}\\[0.2cm]
  \begin{tabular}{@{}cccccc@{}}
    \toprule
    Step & Gradient $g$ & Residual $e$ & $g + e$ & Sent $v$ & New $e = g + e - v$ \\
    \midrule
    1 & 0.4 & 0 & 0.4 & 0 & 0.4 \\
    \uncover<2->{2} & \uncover<2->{0.3} & \uncover<2->{0.4} & \uncover<2->{0.7} & \uncover<2->{\textbf{1}} & \uncover<2->{$-0.3$} \\
    \uncover<2->{3} & \uncover<2->{0.2} & \uncover<2->{$-0.3$} & \uncover<2->{$-0.1$} & \uncover<2->{0} & \uncover<2->{$-0.1$} \\
    \uncover<2->{4} & \uncover<2->{0.4} & \uncover<2->{$-0.1$} & \uncover<2->{0.3} & \uncover<2->{0} & \uncover<2->{0.3} \\
    \uncover<2->{5} & \uncover<2->{0.3} & \uncover<2->{0.3} & \uncover<2->{0.6} & \uncover<2->{\textbf{1}} & \uncover<2->{$-0.4$} \\
    \bottomrule
  \end{tabular}\\[0.3cm]
  {\normalsize\uncover<3->{Without the residual, 0 of a true 1.6 arrives.\quad With it, 2 is sent and $-0.4$ kept: $2 + (-0.4) = 1.6$.}}
  \mlsyscredit{Collective Communication, Example 5.2 and tables 10 and 11.}
  \note{[2.5 min] [CORE] Example 5.2. The compressor rounds to the nearest whole unit, so any value below 0.5 is sent as zero. Without a residual, all five gradients (0.4, 0.3, 0.2, 0.4, 0.3) round to zero, so the model never sees a true total of 1.6. Step 1: with error feedback, the worker keeps what it did not send. Step 2: at step 2 the residual pushes 0.7 over the threshold and 1 is sent; at step 5 again. Step 3: sent plus residual equals the true sum, every step. Summed over steps, transmitted updates equal the true gradients plus the first residual minus the last, so if the residual stays bounded the averages agree. ENGAGE: the gradient is 0.3 every step. How often is a 1 sent? (Expected: about three times in ten steps.) WARN: error feedback is not automatic; the system must store and re-add the residual. FLEX: [CORE].}
\end{frame}

\begin{frame}{1-bit Adam cut communication volume up to 5$\times$}
  \centering
  \safeimg[width=0.92\textwidth,height=3.0cm,keepaspectratio]{onebit.pdf}\\[0.2cm]
  \begin{columns}[c]
    \begin{column}{0.5\textwidth}
      \small Full-precision Adam until the variance settles; then send 1-bit momentum with a local residual.
    \end{column}
    \begin{column}{0.46\textwidth}
      \footnotesize\ttfamily
      "optimizer": \{\\
      \hspace*{1em}"type": "OneBitAdam",\\
      \hspace*{1em}"params": \{ "lr": 4e-4,\\
      \hspace*{2em}\textcolor{commamber}{\textbf{"freeze\_step": 23000}},\\
      \hspace*{2em}"comm\_backend\_name": "nccl" \}\}
    \end{column}
  \end{columns}
  \mlsyscredit{Tang et al. 2021 (1-bit Adam), experiments up to 256 GPUs, via Collective Communication, \S5.9.4. Config: DeepSpeed v0.14.0, docs/\_tutorials/onebit-adam.md (abridged).}
  \note{[1 min] [OPTIONAL] These numbers are measured. Adam divides momentum by the square root of a running variance; noise from 1-bit gradients distorts that denominator. DeepSpeed's team observed the variance settles early, so 1-bit Adam runs normal Adam for a warmup, freezes the variance, then sends only the sign of the momentum plus a scale, keeping the residual locally. Tang et al. report up to 5 times less communication volume (less than the raw 32 times, because warmup and other traffic are counted), up to 3.3 times higher throughput for BERT-Large pretraining and 2.9 times for SQuAD fine-tuning, with convergence matching uncompressed Adam. The config is DeepSpeed's tutorial example, in which freeze\_step is the warmup, 23,000 steps, after which compression begins. ENGAGE: why only 5 times, when 1 bit is 32 times smaller than FP32? (Expected: the uncompressed warmup and the rest of the step's traffic.) LINK: Lecture 11 applies update compression to phone uplinks. FLEX: cut if short.}
\end{frame}

\begin{frame}{Overlap hides each reduction under the next layer's backward}
  \centering
  \mlsyssteps[width=0.95\textwidth,height=5.0cm,keepaspectratio]{overlap_step}{3}\\[0.05cm]
  {\small Exposed communication is the transfer time the step still waits for after overlap.}
  \mlsyscredit{Collective Communication, Napkin Math 5.9 (7B model, 32 layers, 64 GPUs, 100 MB buckets at an assumed 3 ms each) and \S5.10.4. 487.6 ms derived from 0.85 ms per bucket.}
  \note{[3 min] [CORE] Backward runs from the output layer to the input layer, so a layer's gradient is ready while earlier layers still compute. Step 1: sequential. 32 layers of 15 ms of backward, then 32 reductions of 27 ms: 1,344 ms. Step 2: overlapped. Each reduction starts when its layer finishes. But 27 ms is longer than 15 ms, so a backlog of 12 ms grows each layer, and the last reduction has no compute left to hide under: 480 + 27 + 31 times 12 = 879 ms, 34.6 percent less. Step 3: at the hierarchical schedule's rate, about 0.85 ms per 100 MB bucket, a layer's reductions take 7.6 ms, under the 15 ms window, so all but the last hide: 487.6 ms (derived). ENGAGE: what is always exposed? (Expected: the last bucket's transfer.) LINK: this is the Fleet Law's overlap term, $T_{\text{overlap}}$. FLEX: [CORE].}
\end{frame}

\begin{frame}{Runtimes bucket gradients and trace every AllReduce}
  \begin{columns}[T]
    \begin{column}{0.44\textwidth}
      \footnotesize\ttfamily
      \begin{tabular}{@{}l@{}}
        if bucket\_cap\_mb is None:\\
        \hspace*{1em}\# default case\\
        \hspace*{1em}\# (bucket cap is 25 MiB)\\
        \hspace*{1em}\textcolor{commamber}{\textbf{bucket\_cap\_mb = 25}}\\
        ...\\
        self.bucket\_bytes\_cap = int(\\
        \hspace*{1em}bucket\_cap\_mb * 1024 * 1024)
      \end{tabular}\\[0.25cm]
      \normalfont\small Bigger buckets pay fewer startups;\\smaller ones start transfers sooner.
    \end{column}
    \begin{column}{0.54\textwidth}
      \centering
      \safeimg[width=\textwidth,height=4.6cm,keepaspectratio]{horovod_trace.pdf}\\[0.05cm]
      {\footnotesize Horovod Timeline: bytes move for 19 of 325 ms.}
    \end{column}
  \end{columns}
  \mlsyscredit{Code: PyTorch v2.4.0, torch/nn/parallel/distributed.py, lines 793 to 799 (comment line split for the slide). Trace: Sergeev and Del Balso 2018, Figure 5, redrawn (phase edges read to about 2 ms).}
  \note{[2 min] [OPTIONAL] Left: real PyTorch source. DistributedDataParallel groups gradients into buckets, 25 MiB by default, and launches one asynchronous AllReduce per bucket from gradient hooks as soon as all its gradients are ready. Larger buckets amortize per-call startup; smaller ones start earlier and widen the overlap window. One late gradient holds its whole bucket. Right: a real trace from Horovod's timeline tool, redrawn from the paper's screenshot. Each row is one gradient tensor's AllReduce across all ranks. Violet is negotiation, the wait until every rank reports that tensor ready; gray is waiting for the gradient data on the GPU, its own or that of tensors fused with it; amber is NCCL moving bytes. Tensors 0 and 1 finish in 115 ms. Tensors 2 to 4 negotiate for 131 ms, and of the 325 ms from first request to finish, NCCL moves bytes for 19, about 6 percent (derived in calcs.py; the 184 ms AllReduce span is printed in the screenshot). WARN: async\_op=True only returns control to Python; LogP separates network time, which compute can hide, from the processor overhead of starting and finishing each transfer, which it cannot. Calibrate $\alpha$ and $\beta$ with nccl-tests on the job's own placement; Lecture 9 attributes step time. FLEX: cut if short.}
\end{frame}

% =============================================================================
% SECTION: WRAP-UP
% =============================================================================
\mlsysnote{[0.5 min] [CORE] Check the misconceptions, then return to Horovod.}
\mlsysdivider{What decides how long N accelerators take to agree?}{Collective Communication, \S5.12 and \S5.13}
\section{Wrap-Up}

\begin{frame}{True or false?}
  \mlsysactivity{Vote}{2 min}
  \small
  \renewcommand{\arraystretch}{1.35}
  \begin{tabular}{@{}>{\raggedright\arraybackslash}m{0.44\textwidth}>{\raggedright\arraybackslash}m{0.50\textwidth}@{}}
    \toprule
    \textbf{Claim} & \uncover<2->{\textbf{Verdict}} \\
    \midrule
    Bandwidth is the only metric that matters. & \uncover<2->{\textbf{\textcolor{RedLine}{False.}} A 24 KB ring chunk is latency-bound.} \\
    Compression is safe without error feedback. & \uncover<2->{\textbf{\textcolor{RedLine}{False.}} Dropped values never arrive: 0 of 1.6.} \\
    An async collective always hides its latency. & \uncover<2->{\textbf{\textcolor{RedLine}{False.}} Overhead and the last bucket stay exposed.} \\
    A flat ring is good enough across nodes. & \uncover<2->{\textbf{\textcolor{RedLine}{False.}} It runs at one port: 40 ms against 8.3.} \\
    \bottomrule
  \end{tabular}
  \mlsyscredit{Collective Communication, \S5.12 (Fallacies and Pitfalls).}
  \note{[2 min] [CORE] Students vote on each claim, then reveal. Each verdict carries the lecture's number. Bandwidth: Napkin Math 5.1's 24.4 KB chunks on 1,024 ranks are below the 75 KB knee, so more bandwidth helps little. Compression: Example 5.2's rounding sends nothing at all without a residual. Async: Napkin Math 5.9 leaves the last transfer exposed, and LogP's processor overhead cannot be hidden. Flat ring: Napkin Math 5.4. ENGAGE: which claim did your neighbor believe at the start? FLEX: [CORE]. An extension, if time remains: 'rank-to-GPU mapping does not affect collective performance' (false; misplaced rails give one eighth the injection bandwidth).}
\end{frame}

\begin{frame}{Each Horovod measurement moved one term}
  \begin{columns}[c]
    \begin{column}{0.40\textwidth}
      \centering\safeimg[width=\textwidth,height=3.0cm,keepaspectratio]{horovod-timeline.png}\\[0.05cm]
      {\footnotesize Horovod's timeline exposed many tiny AllReduces in ResNet-101.}
    \end{column}
    \begin{column}{0.58\textwidth}
      \footnotesize
      \renewcommand{\arraystretch}{1.15}
      \begin{tabular}{@{}>{\raggedright\arraybackslash}p{3.9cm}>{\raggedright\arraybackslash}p{4.3cm}@{}}
        \toprule
        Horovod measured & Term it moved \\
        \midrule
        Ring for parameter server: about 50\% $\rightarrow$ 88\% & bytes no longer concentrate at one port \\
        \uncover<2->{Tensor Fusion: up to 65\%} & \uncover<2->{startup $\alpha$, paid once per fused buffer} \\
        \uncover<3->{RDMA: +3--4\% ResNet, +30\% VGG-16} & \uncover<3->{per-message cost; VGG-16 is byte-heavy} \\
        \bottomrule
      \end{tabular}
    \end{column}
  \end{columns}
  \mlsyscredit{Sergeev and Del Balso 2018, Sections 7 and 8 (128 Pascal GPUs; Tensor Fusion on TCP; RDMA against 25GbE TCP) and Figure 5 (timeline).}
  \note{[2.5 min] [CORE] Return to the hook. Step 1: replacing the parameter server with a ring removed the concentration that sends $N$ times the bytes through one port, as on our scoreboard (1,280 ms against 39.4 ms, and 8.3 ms with hierarchy); efficiency went from about half to 88 percent. Step 2: Horovod's timeline showed ResNet-101 issuing many tiny AllReduces; Tensor Fusion packs ready tensors into a 64 MB buffer, the fusion lever for startup-bound messages, and gave up to 65 percent on models with many layers over TCP. Step 3: RDMA, kernel-bypass networking, added only 3 to 4 percent for Inception V3 and ResNet-101 but 30 percent for VGG-16. The paper explains that VGG-16's many parameters and few layers shift its critical path from GPU computation to communication, so only there does a cheaper transfer shorten the step. ENGAGE: why did RDMA help one model ten times more than the others? (Expected: only VGG-16 had communication on its critical path.) FLEX: [CORE].}
\end{frame}

\begin{frame}{Takeaways}
  \normalsize
  \small
  \begin{itemize}\setlength{\itemsep}{0.15cm}
    \item Choose the narrowest contract the next computation consumes.
    \item Each round pays $\alpha$ on a chunk of $M/N$; the ring sends the minimum bytes in $2(N-1)$ rounds.
    \item Trees cut rounds; hierarchy shrinks the slow tier's payload to $M/G$ when each GPU has its own port.
    \item An AllToAll finishes when its hottest receiver or its narrowest cut finishes.
    \item Compression and overlap act on what remains; judge them by step time and time to quality.
  \end{itemize}
  \note{[1 min] [CORE] One line per segment. Ask students to star the one they would have gotten wrong before today.}
\end{frame}

\begin{frame}{Open questions}
  \normalsize
  \begin{itemize}\setlength{\itemsep}{0.35cm}
    \item In-switch reduction moves arithmetic into the network. Which other parts of a training step belong there?
    \item Routing skew is set by the data. Should the router balance load, or should the fabric absorb skew?
    \item Overlap needs memory for in-flight buffers. When does hiding communication cost more memory than it saves time?
  \end{itemize}
  \note{[0.5 min] [OPTIONAL] No single answer. FLEX: use one as a two-minute discussion if time remains, or as a written assignment. For the first, point students to SHARP's limits in \S5.7.3: supported operators and formats, and a finite number of aggregation trees per switch. For the second, the auxiliary load-balancing loss changes the model to fit the network; Lecture 6 returns to it.}
\end{frame}

\begin{frame}{Read and explore}
  \small
  \begin{itemize}\setlength{\itemsep}{0.08cm}
    \item \textbf{Read:} Collective Communication, \S5.3 to \S5.10 and the self-checks. \textbf{Practice:} redo Napkin Math 5.3 for 10 MB; run nccl-tests on two GPUs and fit $\alpha$, $\beta$.
    \item Sergeev and Del Balso, ``Horovod: Fast and Easy Distributed Deep Learning in TensorFlow,'' 2018.
    \item Patarasuk and Yuan, ``Bandwidth Optimal All-reduce Algorithms for Clusters of Workstations,'' JPDC 2009.
    \item Karimireddy et al., ``Error Feedback Fixes SignSGD and Other Gradient Compression Schemes,'' ICML 2019.
    \item Li et al., ``PyTorch Distributed: Experiences on Accelerating Data Parallel Training,'' VLDB 2020.
  \end{itemize}
  \note{[0.5 min] [CORE] The chapter first; its napkin math repeats today's numbers. Horovod is the hook's source; Patarasuk and Yuan prove the ring's byte bound; Karimireddy et al. is the error-feedback analysis; Li et al. describe DDP's buckets and overlap from the PyTorch team.}
\end{frame}

\begin{frame}[c]
  \centering
  {\small\textcolor{lightgray}{NEXT LECTURE \textbar{} DISTRIBUTED TRAINING}}\\[0.5cm]
  {\LARGE\bfseries How should a model too large for one node\\be split across a thousand accelerators?}\\[0.3cm]
  \mlsysphoto[height=2.6cm,keepaspectratio]{google-tpu-v4-board.jpg}\\[0.15cm]
  {\normalsize\textcolor{midgray}{Today priced every exchange. Next, choose the split that creates them.}}
  \mlsyscredit{Photo: a TPU v4 board; Norman P. Jouppi et al., CC BY 4.0, Wikimedia Commons.}
  \note{[0.5 min] [CORE] Today took the traffic as given and priced it. Lecture 6 chooses the split: data, sharded, tensor, pipeline, and expert parallelism, and uses today's prices to judge when each repays the communication it creates. The photo is a Google TPU v4 board; pods of such chips reduce along the axes of a torus, the schedule in \S5.7.2.}
\end{frame}

\end{document}
